% ID: SEL-D0010-BD-E07
% ESTADO: cotejada. Segunda pasada visual completa y borrador compilado revisados.
% Fuente: Selectividad/matematicas_aplicadas/EXTRA_TITULAR_MAT-APLIC-CCSS-JULIO-IMPRESO.pdf
% SHA-256: 548f60ee334516c7fa918dba3c28fb3ea5fe753cf76eb05402d7b8acb5a3b560
% Página: 2; bloque D; ejercicio original 7.

Para estimar la proporción de residentes británicos en España que están
a favor de la salida del Reino Unido de la Unión Europea (UE), se toma una
muestra aleatoria de 250 de estos residentes, obteniéndose que 115 estaban
a favor de dejar de pertenecer a la UE.
\begin{enumerate}[label=\alph*)]
  \item \textbf{(1.5 puntos)} Calcule un intervalo de confianza al 99.5\%,
  para estimar la proporción real de esos residentes que está a favor de la
  salida del Reino Unido de la UE.
  \item \textbf{(1 punto)} Manteniendo la misma proporción muestral y el
  mismo nivel de confianza del apartado anterior, determine el tamaño mínimo
  necesario de la muestra, para estimar la proporción de residentes británicos
  en España que están a favor de la salida del Reino Unido de la UE, con un
  error inferior al 5\%.
\end{enumerate}
